\documentclass{article}
\usepackage[margin=1in]{geometry}
\usepackage{graphicx}
\usepackage{amsmath}
\usepackage{booktabs}
\usepackage{natbib}
\usepackage{pgfplotstable}
\usepackage{placeins}
\usepackage[hypertexnames=false]{hyperref}
\renewcommand{\sectionautorefname}{Section}
\renewcommand{\subsectionautorefname}{Section}
\renewcommand{\subsubsectionautorefname}{Section}

\begin{document}
\title{Subgrid and Cosmological parameter variations}
\maketitle
\section{Introduction}

We have expanded the simulation--emulator--calibration framework to a combined cosmology and subgrid parameter regime. Subgrid effects depend on environmental factors like local density, which are themselves strongly influenced by cosmological parameters such as $\sigma_8$. It is therefore important to understand the interactions and correlations between subgrid and cosmology parameters, their impact on the forward models, and the resulting systematics on the calibrations.

Related work: CAMELS with cosmological back-reaction of baryons, \url{https://arxiv.org/pdf/2601.06258}.

The full set of figures supporting the three classes of results sit in three companion notes:
\begin{itemize}
  \item \texttt{summaries\_emulation.tex} -- emulation of summary statistics: $P(k)$ suppression, HMF, GSMF, $f_{\rm gas}$, CSFR.
  \item \texttt{cluster\_profiles.tex} -- emulation of cluster radial profiles (CGD, CGED, CEP, CEEP, CMP, CPP, CTP, CYP).
  \item \texttt{calibration\_results.tex} -- MCMC calibration results and the v1$\leftrightarrow$v2 pipeline comparison.
\end{itemize}
This document collects the simulation-setup details and parameter tables that the other three notes refer to but do not reproduce.

\section{Simulation Framework}
\label{sec:data}


There are five subgrid parameters
$\theta_{\rm sub}=\{\kappa_\mathrm{w}, e_\mathrm{w}, M_\mathrm{seed}, v_\mathrm{kin}, \epsilon_\mathrm{kin}\}$,
and two $\Lambda$CDM parameters
$\theta_{\rm CDM}=\{\omega_m,\sigma_8\}$.


All other parameters in the full subgrid prescription are held fixed (see \cite{frontiere2025modeling} for the parameter choices). The remaining $\Lambda$CDM parameters are also fixed:
$\Omega_{\rm b}h^2=0.02242$, $H_0=67.66$, $n_s=0.9665$. Spatial flatness and massless neutrinos are assumed.

The simulation campaign comprises 110 HAvoCC-400\,Mpc runs assembled from four sub-designs:
\begin{itemize}
    \item indices 0--39: 40 D-optimal points spanning the ranges of \autoref{tab:param_ranges},
    \item indices 40--89: 50 space-filling points over the same box,
    \item indices 90--99: 10 additional space-filling points,
    \item indices 100--110: 10 further space-filling points.
\end{itemize}

The exhaustive parameter listing is in \autoref{appendix:params}; design-space projection plots (\texttt{exp\_design.png}) live in \texttt{summaries\_emulation.tex}. Five simulations, indices $\{3,11,19,27,35\}$, are held out as the validation set used throughout the companion notes.

\begin{table}[htbp]
\centering
\caption{Subgrid and cosmology parameter ranges explored in the simulation campaign.}
\label{tab:param_ranges}
\begin{tabular}{lc}
\hline\hline
Parameter & Range \\
\hline \\
\multicolumn{2}{l}{\textbf{Subgrid Parameter Ranges}} \\
$\kappa_w$ & $2 \le \kappa_w \le 4$ \\
$e_w$ & $0.2 \le e_w \le 1$ \\
$M_\mathrm{seed}$ & $0.6 \times 10^{6} \le M_\mathrm{seed} \le 2.0 \times 10^{6}$ \\
$v_\mathrm{kin}$ & $0.1 \times 10^{4} \le v_\mathrm{kin} \le 1.2 \times 10^{4}$ \\
$\epsilon_\mathrm{kin}$ & $0.02 \times 10^{1} \le \epsilon_\mathrm{kin} \le 1.2 \times 10^{1}$ \\ \\
\hline \\
\multicolumn{2}{l}{\textbf{Cosmology Parameter Ranges}} \\
$\omega_m$ & $0.12 \le \omega_m \le 0.155$ \\
$\sigma_8$ & $0.7 \le \sigma_8 \le 0.9$ \\
\hline\hline
\end{tabular}
\end{table}


\subsection{Statistics for emulation and calibration}
\label{sec:obs}


The emulator statistics considered here are
\begin{enumerate}
  \item Galaxy stellar mass function (GSMF),
  \item Halo mass function (HMF),
  \item Cluster gas fraction $f_\mathrm{gas}$,
  \item Total matter power-spectrum suppression $P_\mathrm{hydro}(k)/P_\mathrm{GO}(k)$,
  \item Cosmic star formation rate (CSFR),
  \item Eight cluster radial profiles: gas density (CGD), gas electron density (CGED), gas pressure (CPP), gas temperature (CTP), gas entropy (CEP), electron entropy (CEEP), gas metallicity (CMP) and Compton-$y$ (CYP).
\end{enumerate}
See \texttt{summaries\_emulation.tex} and \texttt{cluster\_profiles.tex} for training-ensemble and validation figures.

\subsection{Power spectra suppression}
We compute the ratio of the total matter power spectra from a hydrodynamical simulation and its gravity-only counterpart:

\begin{equation}
\frac{P_{\rm hydro}\!\left(k;\,\theta_{\rm sub}, \theta_{\rm cos}\right)}
     {P_{\rm GO}\!\left(k;\,\theta_{\rm cos}\right)} .
\end{equation}

Here, $\theta_{\rm sub}$ and $\theta_{\rm cos}$ are subgrid parameters and cosmological parameters respectively. The gravity-only (GO) power spectrum in the denominator can be approximated as:

\begin{equation}
P_{\rm GO}\!\left(k;\,\theta_{\rm cos}\right)
\approx
P_{\rm GO}\!\left(k;\,\theta_{\rm cos}'\right)
\frac{P_{\rm emu}\!\left(k;\,\theta_{\rm cos}\right)}
     {P_{\rm emu}\!\left(k;\,\theta_{\rm cos}'\right)} ,
\end{equation}
where $P_{\rm GO}\!\left(k;\,\theta_{\rm cos}'\right)$ is the single gravity-only simulation that has been run at the cosmological parameter set $\theta_{\rm cos}'$. $P_{\rm emu}\!\left(k;\,\theta_{\rm cos}\right)$ is output from Cosmic emulator (based on Mira-Titan Universe).

\section{Emulation framework}

The emulator development in this work is carried out with SEPIA (Simulation-Enabled Prediction, Inference, and Analysis), a Python code developed at Los Alamos National Laboratory \citep{james_gattiker_2020_4048801} that implements the Bayesian emulation and calibration methodology described in \cite{Higdon2008}.

For each summary statistic we train an independent multivariate Gaussian-process emulator over the seven-parameter input space. The simulation output is projected onto a principal-component basis (training-time \texttt{exp\_variance}$=0.999$ for all summaries except $P(k)$, which uses $0.95$); the GP is then fit on the PC weights. At inference time the basis size is auto-synced to whatever was stored in the model pickle so chains and emulators stay consistent across retrains.

\emph{Train/test split.} The five held-out simulations $\{3,11,19,27,35\}$ form a fixed validation set; the remaining sims train the per-observable GPs. Per-observable validation, sensitivity, and redshift sweeps are collected in \texttt{summaries\_emulation.tex} ($P(k)$, HMF, GSMF, $f_\mathrm{gas}$, CSFR) and \texttt{cluster\_profiles.tex} (eight cluster profiles).

\section{Redshift coverage and snapshot windows}\label{sec:z_dependence}

Multi-snapshot emulators are trained per-observable on the cleanest redshift window for that statistic. Based on direct inspection of the simulation outputs we restrict each observable to the redshift interval over which the data are clean and sufficiently complete. The snapshot ranges (in the indexing used by our pipeline) and the motivation for each cut are:

\begin{itemize}
  \item \textbf{GSMF, HMF:} snapshots 205--624 (start index 0), corresponding to $z \in [0,\,2.0]$. All measurements across this range are clean.
  \item \textbf{$f_{\mathrm{gas}}$:} snapshots 310--624 (start index 4), $z \in [0,\,1.0]$. Earlier snapshots have bad values: snapshot 205 has 20/39 simulations flagged as bad, and snapshots 224--275 show intermittent issues (typically 1--2 problematic simulations per snapshot).
  \item \textbf{Cluster profiles (CGD, CGED, CEP, CEEP, CMP, CPP, CTP, CYP):} snapshots 415--624 (start index 6), $z \in [0,\,0.5]$. Snapshots 205--247 are entirely NaN and snapshots 275--355 contain partial NaNs.
  \item \textbf{$P(k)$ suppression:} $z=0$ only.
  \item \textbf{CSFR:} continuous in scale factor $a\in[0,1]$.
\end{itemize}

With the available snapshots we first fit independent GP models for each redshift, then a redshift interpolator stitches them together at inference time (used by the multi-$z$ likelihood).


\section{Parameter Constraints Using MCMC}\label{sec:mcmc}

The full MCMC trial set, posteriors and v1$\leftrightarrow$v2 benchmarking are documented in \texttt{calibration\_results.tex}. The shared configuration (priors, walker count, fixed cosmology when applicable) is held in \texttt{Inference/configs/\_defaults.yaml}; trial-specific overrides specify only the observables and the parameter mode (\texttt{subgrid}, \texttt{subgrid+cosmo}, or \texttt{custom}).

\section{Discussion and outlook}
\label{sec:conclusions}

The framework currently delivers (i) single- and multi-$z$ emulators for $P(k)$ suppression, HMF, GSMF, $f_{\rm gas}$, CSFR and the eight cluster profiles, (ii) a redshift interpolator that stitches per-snapshot GPs at inference time, and (iii) a YAML-driven MCMC pipeline that supports arbitrary observable + parameter combinations, including the cluster-only and v1$\leftrightarrow$v2 benchmark fits documented in \texttt{calibration\_results.tex}.

\paragraph{Open items.}
\begin{enumerate}
    \item \textbf{Multi-$z$ GSMF MCMC.} The per-snapshot likelihood scaffolding is in place; the remaining work is a trial config wired to multiple Driver+2022 redshift bins and a follow-up sensitivity test on $\sigma_8$--$M_{\rm seed}$ degeneracy as a function of $z$.
    \item \textbf{Joint $P(k)$ + cluster constraints.} Combine the validated $P(k)$-suppression emulator with $f_{\rm gas}$ (and optionally CGD) for an independent baryonic-feedback channel; check whether the GSMF-driven posterior shifts.
    \item \textbf{Cluster-profile observational coverage.} CEEP and CYP currently have no external dataset in the comparison figure; resolve whether to add literature values or label them prediction-only. Address mild small-$r$ tension in the metallicity (CMP) profile and the outer-bin tension in the temperature (CTP) profile.
    \item \textbf{Design-extension review.} After the next batch of simulations, re-check the design-space coverage plot and decide whether to expand the $\epsilon_{\rm kin}$ range (its posterior currently sits at the prior edge in the 7p fit).
    \item \textbf{Bias-parameter MCMC.} The \texttt{GSMF\_CGD\_bias} trial folds in three observational bias parameters ($\log b_\star, b_{\rm CV}, b_{\rm HSE}$); produce a paper-grade posterior with the same priors as the headline 7p fit.
\end{enumerate}


\section{Parameters in the design}
\label{appendix:params}
\newcommand{\DesignTableSlice}[4]{%
\begin{table}
\centering
\caption{#4}
\label{#3}
\pgfplotstabletypeset[
  col sep=comma,
  columns/kappa_W/.style={column name=$\kappa_w$},
  columns/e_W/.style={column name=$e_w$},
  columns/M_seed/.style={column name=$M_\mathrm{seed}$},
  columns/v_kin/.style={column name=$v_\mathrm{kin}$},
  columns/eps_kin/.style={column name=$\epsilon_\mathrm{kin}$},
  columns/omega_m/.style={column name=$\omega_m$},
  columns/sigma_8/.style={column name=$\sigma_8$},
  every head row/.style={before row=\toprule, after row=\midrule},
  every last row/.style={after row=\bottomrule},
  row predicate/.code={%
    \ifnum\pgfplotstablerow<#1\relax\pgfplotstableuserowfalse\fi
    \ifnum\pgfplotstablerow>#2\relax\pgfplotstableuserowfalse\fi
  },
]{FinalDesign.txt}
\end{table}
}

\DesignTableSlice{0}{39}{tab:params_dopt}{Design points 0--39: 40 simulations of D-optimal design around \autoref{tab:param_ranges}.}

\DesignTableSlice{40}{89}{tab:params_spacefill_A}{Design points 40--89: 50 space-filling design around \autoref{tab:param_ranges}.}

\DesignTableSlice{90}{99}{tab:params_spacefill_B}{Design points 90--99: 10 space-filling design.}

\DesignTableSlice{100}{110}{tab:params_spacefill_C}{Design points 100--110: 10 additional space-filling design points.}

\clearpage

\bibliographystyle{aasjournal}
\bibliography{biblio}
\end{document}
